\paragraph{} The open circuit and short circuit tests reveal that in real transformers there are
measurable non-idealities. Being able to derive the values of $R_{c}$, $X_{m}$, $R_{eq}$, and $X_{eq}$
is highly desirable in real life applications. While these values are relatively small in our lab tests,
they would easily become much larger and more significant in practical situations. It's not hard to imagine
that in a transformer responsible for stepping down transmission line level voltage (~700kV) to a
sub-transmission level voltage, that these values could become quite large and crucial to determine
in order to manage efficiency and safety. Similarly, we are sure that these sorts of tests are applicable to
other areas of electrical engineering. Often there is not a perfect way to calculate non-idealities,
and deriving them experimentally is the best way. The general method of performing this test that the
lab introduced us to is certain to be pertinent in other topics we encounter in the future.